%=============================================================================!
% 5.0  CHAPTER OPENING                                                        !
%=============================================================================!
Chapters~2 to 4 treated a body as a particle when its motion could be
specified by a single position. A molecule also changes its orientation
as it moves, so we now need to describe rotation and distinguish it from
motion of the centre of mass. Angular momentum supplies the corresponding
conservation law: for an isolated group with central pair forces, the
internal torques cancel and the total angular momentum stays fixed.

We begin with a puck moving in a plane, where the geometry can be followed
through its distance and angle from a fixed point. Extending that argument
to space introduces the cross product; applying it to many particles then
leads to the inertia tensor, which relates angular velocity to angular
momentum. These tools let us separate a molecule's drift and rotation
from its internal motion, and explain the count of $3N - 6$ vibrational
degrees of freedom, or $3N - 5$ for a linear molecule. Chapter~6 uses the
same distinction between allowed motions to choose coordinates for
constrained systems.
